\honest{Professing and Cheering}{Eliezer Yudkowsky}{2007}

I once attended a panel on whether science and religion are compatible. One panelist, a pagan, spent what felt like forever telling how the Earth was made: a primordial cow licked a god into existence, and the god's descendants killed a giant and made the Earth from his corpse. \nb{This is the Norse creation story as Snorri Sturluson tells it in the \textsc{Prose Edda}.} It was more absurd than a world on a turtle's back, and she ``clearly knew enough science to know this.''

I still struggle to name what I saw: ``pride? Self-satisfaction? A deliberate flaunting of herself?'' She said religion and science were compatible. She even explained that the Vikings talked of a primordial abyss because of the land they lived in, and yet she insisted that this was what she ``believed.'' \nb{Both remarks fit a reading the post does not consider: that she took the myth as symbolic. Wikipedia's article on modern Heathenry, citing academic studies, says most practitioners take the Norse cosmology as ``a poetic or symbolic description of the cosmos,'' while believing that the gods themselves exist. The post does not report asking her what she meant by ``believed.''}

Dennett's ``belief in belief'' does not fit her, I say. She showed no fanatical faith and did not need us to believe her. Dennett has also suggested studying religious belief as ``religious profession,'' the way an anthropologist might ask why students all write the same strange sentence on quizzes. \nb{This could not be checked against Dennett's book.} I think Dennett is slightly too cynical, because ``most people are honest enough'' to repeat to themselves what they say aloud. I give no evidence for this.

Profession does not fit her either. Someone professing a belief to satisfy a priest, a fellow believer or their own self-image would pretend much more convincingly, and she was not even trying.

So it ``finally occurred to me'' that she was ``cheering for paganism,'' like holding up a GO BLUES banner, which states no fact and persuades no one. Her pride was like marching naked in a gay pride parade: an outrageous cheer, made in the belief that she ``couldn't be arrested or criticized.'' That is why the absurdity mattered to her. A plausible version would have been ``like putting on clothes.'' \nb{What the post says she knew, intended and believed rests on the author's impression of her manner. It cites nothing she said in support.}

In a footnote I add that there is nothing wrong with actually marching naked in pride parades.
